\section*{Bab \ref{ch:g4:world-physical-map} --- Peta Fisik Dunia}

\begin{solution}{exo:g4:world-physical-map:1}
\omterm{def:g4:world-physical-map:altitude}{Ketinggian} suatu tempat adalah tingginya di atas \omterm{def:g4:world-physical-map:altitude}{permukaan laut}, dalam
meter. \omterm{def:g4:world-physical-map:altitude}{Ketinggian} diukur dari \omterm{def:g4:world-physical-map:altitude}{permukaan laut}, yaitu \omterm{def:g4:world-physical-map:altitude}{ketinggian} rata-rata
permukaan air \omterm{def:g2:water-on-earth:ocean}{laut}.
\end{solution}

\begin{solution}{exo:g4:world-physical-map:2}
(a) Hijau tua: antara 0 dan 200 meter. (b) Kuning pucat: antara 500 dan
1\,000 meter. (c) Hampir putih: lebih dari 4\,000 meter.
\end{solution}

\begin{solution}{exo:g4:world-physical-map:3}
(a) Amerika Selatan. (b) Asia. (c) Afrika. (d) Amerika Utara.
\end{solution}

\begin{solution}{exo:g4:world-physical-map:4}
\omterm{def:g1:landscapes-around-us:landform}{Gunung} tertinggi adalah Gunung Everest (8\,849 meter). \omterm{def:g1:landscapes-around-us:water}{Danau} terbesar adalah
\omterm{def:g2:water-on-earth:ocean}{Laut} Kaspia; \omterm{def:g1:landscapes-around-us:water}{danau} terdalam adalah \omterm{def:g1:landscapes-around-us:water}{Danau} Baikal.
\end{solution}

\begin{solution}{exo:g4:world-physical-map:5}
Makin tua warna birunya, makin dalam \omterm{def:g2:water-on-earth:ocean}{lautnya}: \omterm{def:g2:water-on-earth:ocean}{laut} dangkal di dekat \omterm{def:g1:landscapes-around-us:water}{pesisir}
(kurang dari 200 meter) dan jauh lebih dalam di tengah \omterm{def:g2:water-on-earth:ocean}{samudra} (lebih dari
4\,000 meter di banyak tempat).
\end{solution}

\begin{solution}{exo:g4:world-physical-map:6}
\omterm{def:g4:world-physical-map:physical-map}{Peta fisik} menunjukkan \omterm{def:g4:world-physical-map:altitude}{ketinggian} permukaan. Di Antarktika dan Greenland,
permukaannya adalah bagian atas lapisan es yang tebal, yang terletak tinggi
di atas \omterm{def:g4:world-physical-map:altitude}{permukaan laut}, sehingga keduanya mendapat warna daratan tinggi.
\end{solution}

\begin{solution}{exo:g4:world-physical-map:7}
Dasar Challenger Deep (hampir 11\,000 meter di bawah \omterm{def:g4:world-physical-map:altitude}{permukaan laut}); dasar
\omterm{def:g2:water-on-earth:ocean}{samudra} rata-rata (kira-kira 3\,700 meter di bawahnya); sebuah pantai
(kira-kira 0 meter); sebuah desa pada 1\,500 meter; puncak Gunung Everest
(8\,849 meter).
\end{solution}

\begin{solution}{exo:g4:world-physical-map:8}
Sebuah \omterm{def:g1:landscapes-around-us:water}{sungai} terbentuk dari banyak \omterm{def:g1:landscapes-around-us:water}{sungai} kecil yang bergabung, dan orang
tidak sepakat tentang \omterm{def:g1:landscapes-around-us:water}{sungai} kecil mana yang merupakan awalnya yang
sebenarnya. Jika diukur dari sumber yang lebih jauh, \omterm{def:g1:landscapes-around-us:water}{sungai} itu menjadi
lebih panjang. Panjang Nil dan Amazon begitu mirip (masing-masing kira-kira
6\,700 sampai 6\,800 kilometer) sehingga jawabannya bergantung pada pilihan
itu.
\end{solution}

\begin{solution}{exo:g4:world-physical-map:9}
Seperempat dari 30 juta adalah tujuh setengah juta kilometer persegi (30
dibagi 4); sepertiganya adalah 10 juta (30 dibagi 3). Sahara, delapan
setengah juta, lebih dari seperempat Afrika tetapi kurang dari sepertiganya.
\end{solution}

\begin{solution}{exo:g4:world-physical-map:10}
$2500 - 1000 = 1500$: pondok itu 1\,500 meter lebih tinggi. Pondok itu,
antara 2\,000 dan 4\,000 meter, berwarna cokelat; \omterm{def:g1:landscapes-around-us:water}{danau} itu, pada
1\,000 meter, berada di batas antara kuning pucat (500 sampai 1\,000) dan
cokelat muda (1\,000 sampai 2\,000): keduanya diterima.
\end{solution}

\begin{solution}{exo:g4:world-physical-map:11}
10 kilometer sama dengan 10\,000 meter. $10000 - 8849 = 1151$: pesawat itu
lebih tinggi daripada Gunung Everest, selisihnya 1\,151 meter. Di atas
Challenger Deep, pesawat itu berada 10\,000 meter di atas \omterm{def:g2:water-on-earth:ocean}{laut}, dan dasar
palung terletak 11\,000 meter di bawah \omterm{def:g2:water-on-earth:ocean}{laut}: $10000 + 11000 = 21000$ meter
di atas dasar palung.
\end{solution}
